% Section: Abstract
Length-controlled evaluation metrics have been proposed to debias LLM preference assessments,
but whether they preserve the capability ordering that human annotators reveal has not been
tested with confound-controlled methods at population scale. We ask: after explicitly removing
the effect of response verbosity, does human preference (\texttt{win\_rate}) still predict
length-debiased preference (\texttt{LC\_winrate}) --- and how strongly? Analyzing the
AlpacaEval 2.0 leaderboard across 223 diverse LLMs, we find a Spearman partial correlation of
0.985 ($p = 1.69 \times 10^{-170}$, one-tailed), higher than the 0.94 bivariate correlation
previously reported --- capability explains ${\sim}5\times$ more LC variance than verbosity in
standardized regression ($|\beta_{\mathrm{win}}| = 21.34$ vs $|\beta_{\mathrm{len}}| = -4.37$),
and the monotonic capability-LC ordering across quartiles carries a large effect size
($\varepsilon^2 = 0.883$). Two independent statistical estimators converge within 1.1pp
(Frisch--Waugh--Lovell verification), ruling out methodological artifact. These results validate
length-controlled evaluation as a capability-order-preserving transformation, and establish a
concrete lesson for evaluation study design: using the alignment gap $\Delta$ rather than
\texttt{LC\_winrate} as outcome attenuates effect sizes by ${\sim}10\times$ due to mathematical
composition.
